% ============================================================
% SECTION: Skema Database dan Deployment ke VPS
% Akan di-\input ke bab3.tex
% ============================================================

% -------------------------------------------------------
\section{Skema Database dan Manajemen Infrastruktur}
\label{sec:database_infrastruktur}
% -------------------------------------------------------

Bagian ini menjelaskan secara rinci desain skema database PostgreSQL dan seluruh
infrastruktur teknis yang mendukung sistem NusaNara di lingkungan produksi.
Pemahaman skema dan infrastruktur sangat penting karena seluruh pipeline RAG ---
dari penyimpanan vektor hingga sistem adaptif --- bergantung pada keputusan desain
di lapisan ini.

% -------------------------------------------------------
\subsection{Skema Database PostgreSQL}
\label{subsec:db_schema}
% -------------------------------------------------------

Skema database terdiri dari tiga tabel utama yang saling berelasi. Seluruh tabel
berada dalam satu instance PostgreSQL 16, memanfaatkan ekstensi \textbf{pgvector}
untuk kolom bertipe \texttt{VECTOR(768)}.

\subsubsection{Tabel 1: \texttt{knowledge\_base}}

Tabel ini merupakan \textit{Single Source of Truth} untuk seluruh konten lowongan
pekerjaan yang menjadi basis pengetahuan RAG. Setiap baris merepresentasikan satu
dokumen lowongan.

\begin{lstlisting}[language=SQL, caption={DDL Tabel knowledge\_base}]
CREATE EXTENSION IF NOT EXISTS vector;

CREATE TABLE knowledge_base (
    id              SERIAL PRIMARY KEY,
    cluster         TEXT NOT NULL,          -- Klaster karier (8 nilai)
    title           TEXT NOT NULL,          -- Judul posisi lowongan
    company         TEXT NOT NULL,          -- Nama perusahaan
    location        TEXT,                   -- Kota/lokasi kerja
    salary_text     TEXT,                   -- Teks gaji asli, e.g. "2,5 jt - 3,5 jt"
    salary_min      INTEGER,                -- Gaji minimum dalam Rupiah (hasil parsing)
    salary_max      INTEGER,                -- Gaji maksimum dalam Rupiah (hasil parsing)
    skills          TEXT[],                 -- Array keterampilan yang diminta
    content         TEXT NOT NULL,          -- Gabungan semua field (untuk embedding)
    embedding       VECTOR(768),            -- Vektor semantik (nomic-embed)
    search_vector   TSVECTOR,               -- Indeks full-text search
    created_at      TIMESTAMPTZ DEFAULT NOW()
);

-- Indeks IVFFlat untuk pencarian ANN (Approximate Nearest Neighbor)
-- lists = sqrt(720) ~= 27 cluster K-Means
CREATE INDEX idx_knowledge_embedding
    ON knowledge_base
    USING ivfflat (embedding vector_cosine_ops)
    WITH (lists = 27);

-- Indeks GIN untuk pencarian full-text yang cepat
CREATE INDEX idx_knowledge_fts
    ON knowledge_base
    USING GIN (search_vector);

-- Trigger untuk auto-generate search_vector dari kolom content
CREATE OR REPLACE FUNCTION update_search_vector()
RETURNS TRIGGER AS $$
BEGIN
    NEW.search_vector := to_tsvector('simple', NEW.content);
    RETURN NEW;
END;
$$ LANGUAGE plpgsql;

CREATE TRIGGER trig_update_search_vector
    BEFORE INSERT OR UPDATE ON knowledge_base
    FOR EACH ROW EXECUTE FUNCTION update_search_vector();
\end{lstlisting}

\noindent \textbf{Penjelasan desain tabel:}
\begin{itemize}
  \item \textbf{Kolom \texttt{skills} (TEXT[]):} Array PostgreSQL dipilih (bukan tabel
    relasional terpisah) karena skill tidak memerlukan operasi join kompleks ---
    hanya operasi overlap ($\cap$) yang cukup dilakukan di Python setelah fetch.
    Array juga memungkinkan query sederhana: \texttt{WHERE 'Python' = ANY(skills)}.
  \item \textbf{Kolom \texttt{content}:} Berisi gabungan seluruh field dalam satu
    blok teks (format: ``[title] di [company]. Skills: [skills]'').
    Format ini memastikan embedding menangkap konteks penuh dokumen dalam satu vektor.
  \item \textbf{Kolom \texttt{embedding} (VECTOR(768)):} Disimpan sebagai vektor
    768 dimensi. Setiap dimensi merupakan nilai \texttt{float32}. Total storage per
    dokumen untuk vektor: $768 \times 4 \text{ bytes} = 3072 \text{ bytes} \approx 3$ KB.
    Total storage untuk 720 dokumen: $720 \times 3 \text{ KB} \approx 2.16$ MB.
  \item \textbf{Indeks IVFFlat (\texttt{lists = 27}):} Parameter \texttt{lists}
    menentukan jumlah \textit{cluster} K-Means yang dibangun oleh IVFFlat.
    Nilai 27 ditetapkan berdasarkan formula rekomendasi pgvector untuk dataset berukuran sedang:
    $\texttt{lists} = \lceil\sqrt{N}\rceil = \lceil\sqrt{720}\rceil = \lceil 26{,}8 \rceil = 27$.
    IVFFlat membagi 720 vektor ke dalam 27 \textit{cluster}. Saat \textit{query},
    sistem hanya memeriksa cluster terdekat (\texttt{probes} = 1 secara default),
    menghasilkan pencarian \textit{approximate nearest neighbor} (ANN) yang jauh lebih
    cepat dibanding \textit{brute-force} $O(n)$.
  \item \textbf{Trigger \texttt{search\_vector}:} Menggunakan \texttt{to\_tsvector('simple', ...)}
    dengan konfigurasi \texttt{simple} (tanpa stemming bahasa tertentu) agar kata kunci
    teknis seperti ``Python'', ``PostgreSQL'', dan ``UI/UX'' tidak di-stem menjadi
    bentuk lain yang salah.
\end{itemize}

\subsubsection{Tabel 2: \texttt{user\_profiles}}

Tabel ini menyimpan profil adaptif setiap pengguna yang terakumulasi lintas sesi.

\begin{lstlisting}[language=SQL, caption={DDL Tabel user\_profiles}]
CREATE TABLE user_profiles (
    user_id           TEXT PRIMARY KEY,      -- Clerk user ID (eksternal)
    identified_skills TEXT[] DEFAULT '{}',   -- Array skill terakumulasi
    career_interests  TEXT[] DEFAULT '{}',   -- Array minat karier
    profile_summary   TEXT,                  -- Ringkasan narasi terkini
    session_count     INTEGER DEFAULT 0,     -- Jumlah sesi total
    last_active       TIMESTAMPTZ,           -- Timestamp sesi terakhir
    created_at        TIMESTAMPTZ DEFAULT NOW()
);

-- Indeks untuk lookup cepat berdasarkan user_id (sudah PRIMARY KEY)
-- Indeks GIN untuk query array skill
CREATE INDEX idx_user_skills
    ON user_profiles USING GIN (identified_skills);
\end{lstlisting}

\noindent \textbf{Penjelasan desain tabel:}
\begin{itemize}
  \item \textbf{Primary Key \texttt{user\_id}:} Menggunakan Clerk user ID (format:
    \texttt{user\_xxxx...}) sebagai primary key. Tidak menggunakan UUID internal
    karena Clerk user ID sudah unik secara global dan menghindari satu lapisan
    mapping tambahan.
  \item \textbf{Kolom \texttt{identified\_skills} dan \texttt{career\_interests}:}
    Keduanya bertipe \texttt{TEXT[] DEFAULT '\{\}'} (array kosong, bukan NULL).
    Default array kosong memastikan operasi UNION set di Python tidak memerlukan
    pengecekan NULL: \texttt{set(existing\_skills or []) | set(new\_skills)}.
  \item \textbf{Kolom \texttt{session\_count}:} Di-increment setiap kali
    background task adaptive berhasil dieksekusi. Berguna untuk analisis
    longitudinal: seberapa sering pengguna berinteraksi dengan sistem.
  \item \textbf{Indeks GIN pada \texttt{identified\_skills}:} Memungkinkan query
    efisien ``pengguna yang memiliki skill Python'' menggunakan operator
    \texttt{@>} (contains). Berguna untuk fitur analitik di masa depan.
\end{itemize}

\subsubsection{Tabel 3: \texttt{recommendation\_history}}

Tabel ini menyimpan log lengkap setiap sesi rekomendasi sebagai \textit{audit trail}.

\begin{lstlisting}[language=SQL, caption={DDL Tabel recommendation\_history}]
CREATE TABLE recommendation_history (
    id              SERIAL PRIMARY KEY,
    user_id         TEXT NOT NULL REFERENCES user_profiles(user_id),
    session_id      UUID DEFAULT gen_random_uuid(),
    user_narrative  TEXT NOT NULL,         -- Narasi input pengguna
    profile_snapshot JSONB,               -- Snapshot profil saat itu
    retrieved_docs  JSONB,                -- Top-3 dokumen yang digunakan
    llm_response    TEXT,                 -- Respons lengkap LLM
    rag_scores      JSONB,                -- Skor RAG per dokumen
    ttfb_ms         INTEGER,              -- Time to First Byte (ms)
    created_at      TIMESTAMPTZ DEFAULT NOW()
);

-- Indeks untuk riwayat per-user (query ORDER BY created_at)
CREATE INDEX idx_history_user
    ON recommendation_history (user_id, created_at DESC);
\end{lstlisting}

\noindent \textbf{Penjelasan desain tabel:}
\begin{itemize}
  \item \textbf{Kolom \texttt{profile\_snapshot} (JSONB):} Menyimpan \textit{snapshot}
    profil pengguna \textit{pada saat} sesi terjadi, bukan referensi ke profil
    yang bisa berubah. Ini memungkinkan reproduksi hasil (\textit{reproducibility})
    --- mengapa rekomendasi tertentu diberikan pada waktu itu.
  \item \textbf{Kolom \texttt{retrieved\_docs} (JSONB):} Menyimpan Top-3 dokumen
    lengkap (title, company, skills, cosine similarity, skill bonus, final score).
    Data ini digunakan untuk menampilkan transparansi kepada pengguna dan
    untuk evaluasi \textit{offline} kualitas retrieval.
  \item \textbf{Kolom \texttt{rag\_scores} (JSONB):} Menyimpan skor detail:
    \texttt{\{doc\_id, cosine\_sim, skill\_bonus, final\_score, rrf\_score\}}.
  \item \textbf{JSONB vs JSON:} PostgreSQL \texttt{JSONB} dipilih (bukan \texttt{JSON})
    karena JSONB menyimpan data dalam format biner yang ter-parse, memungkinkan
    query langsung ke field nested (\texttt{rag\_scores->>'cosine\_sim'}) dan
    penggunaan indeks GIN pada field JSONB.
  \item \textbf{Kolom \texttt{ttfb\_ms}:} Menyimpan \textit{Time to First Byte}
    dalam milidetik, memungkinkan monitoring performa historis.
\end{itemize}

% -------------------------------------------------------
\subsection{Strategi Deployment: Native VPS Ubuntu dengan systemd}
\label{subsec:deployment_docker}
% -------------------------------------------------------

Sistem NusaNara dideploy ke VPS Ubuntu 22.04 menggunakan strategi \textbf{tanpa Docker}
untuk komponen backend dan database, namun dengan manajemen proses menggunakan
\textbf{systemd} (init system Linux). Keputusan ini diambil atas pertimbangan berikut:

\begin{itemize}
  \item \textbf{Ollama tidak dapat berjalan optimal di dalam container Docker}
    pada VPS tanpa GPU NVIDIA --- Docker memerlukan NVIDIA Container Toolkit
    untuk passthrough GPU, yang tidak relevan pada VPS CPU-only. Ollama berjalan
    langsung di host untuk akses memori yang lebih efisien.
  \item \textbf{PostgreSQL pada Docker mengalami overhead I/O} pada volume mounts.
    Instalasi native PostgreSQL memberikan performa disk I/O yang lebih baik,
    krusial untuk pencarian pgvector yang membaca data vektor berulang kali.
  \item \textbf{Simplifikasi debugging} pada lingkungan produksi: systemd journal
    (\texttt{journalctl}) lebih mudah diakses daripada Docker logs untuk
    monitoring real-time oleh satu developer.
\end{itemize}

\subsubsection{Arsitektur Proses di VPS}

Empat proses berjalan secara bersamaan di VPS, dikelola oleh systemd:

\begin{table}[H]
\centering
\caption{Proses Produksi di VPS dan Manajemen via Systemd}
\label{tab:proses_vps}
\begin{tabular}{llll}
\toprule
\textbf{Proses} & \textbf{Port} & \textbf{Systemd Unit} & \textbf{RAM Estimasi} \\
\midrule
Nginx (Reverse Proxy) & 80, 443 & nginx.service & $\sim$50 MB \\
FastAPI (Uvicorn) & 8000 (internal) & nusanara.service & $\sim$500 MB \\
PostgreSQL 16 & 5432 (internal) & postgresql.service & $\sim$2 GB \\
Ollama (LLM + Embed) & 11434 (internal) & ollama.service & $\sim$8 GB \\
\midrule
\multicolumn{3}{l}{\textbf{Total estimasi}} & \textbf{$\sim$10.5 GB} \\
\bottomrule
\end{tabular}
\end{table}

Sisa $\sim$5.5 GB dari 16 GB RAM digunakan sebagai buffer OS, cache filesystem
PostgreSQL, dan headroom untuk lonjakan penggunaan.

\subsubsection{Konfigurasi Systemd untuk FastAPI}

Systemd digunakan sebagai \textit{process manager} untuk memastikan FastAPI
otomatis restart saat crash dan saat VPS reboot:

\begin{lstlisting}[language=bash, caption={Unit File Systemd untuk FastAPI NusaNara}]
# /etc/systemd/system/nusanara.service
[Unit]
Description=NusaNara FastAPI Backend
After=network.target postgresql.service ollama.service
Requires=postgresql.service

[Service]
Type=simple
User=ubuntu
WorkingDirectory=/home/ubuntu/nusanara/backend
Environment="PATH=/home/ubuntu/.venv/bin"
Environment="PYTHONPATH=/home/ubuntu/nusanara"
EnvironmentFile=/home/ubuntu/nusanara/.env   # Memuat variabel rahasia
ExecStart=/home/ubuntu/.venv/bin/uvicorn \
    main:app \
    --host 127.0.0.1 \
    --port 8000 \
    --workers 1 \
    --loop uvloop \
    --http httptools
Restart=always
RestartSec=5
StandardOutput=journal
StandardError=journal
SyslogIdentifier=nusanara

[Install]
WantedBy=multi-user.target
\end{lstlisting}

\noindent \textbf{Penjelasan konfigurasi systemd:}
\begin{itemize}
  \item \textbf{\texttt{After=postgresql.service ollama.service}:} Memastikan
    PostgreSQL dan Ollama sudah aktif sebelum FastAPI dimulai. Tanpa ini, aplikasi
    bisa startup sebelum database siap, menyebabkan error koneksi saat pertama kali.
  \item \textbf{\texttt{EnvironmentFile=/home/ubuntu/nusanara/.env}:} Memuat
    variabel lingkungan sensitif (DATABASE\_URL, CLERK\_SECRET\_KEY) dari file
    \texttt{.env} yang tidak masuk ke version control (\textit{gitignored}).
    Ini adalah praktik keamanan standar (\textit{12-factor app}).
  \item \textbf{\texttt{--workers 1}:} Menggunakan 1 worker Uvicorn demi 
    menjaga stabilitas dan mencegah isu \textit{out-of-memory} (OOM) mengingat
    sumber daya VPS terbatas (RAM $\sim$16 GB) yang berbagi beban dengan Ollama
    (Llama 3.1) dan PostgreSQL pgvector. Pendekatan 1 worker diasumsikan cukup untuk
    kebutuhan sistem berskala \textit{demo} dalam lingkup penelitian ini tanpa perlu
    menangani pengguna serentak massal.
  \item \textbf{\texttt{--loop uvloop} dan \texttt{--http httptools}:}
    Library C-based yang menggantikan event loop asyncio default Python.
    uvloop 2-4$\times$ lebih cepat dari asyncio default, krusial untuk
    throughput SSE streaming yang tinggi.
  \item \textbf{\texttt{Restart=always} dengan \texttt{RestartSec=5}:}
    Systemd akan restart proses 5 detik setelah crash, memastikan
    \textit{high availability} tanpa intervensi manual.
\end{itemize}

\subsubsection{Konfigurasi Nginx sebagai Reverse Proxy}

Nginx bertindak sebagai gerbang (\textit{gateway}) antara internet dan
layanan internal (FastAPI pada port 8000):

\begin{lstlisting}[language=bash, caption={Konfigurasi Nginx Lengkap untuk NusaNara}]
# /etc/nginx/sites-available/nusanara
server {
    listen 80;
    server_name api.nusanara.app;
    return 301 https://$host$request_uri;  # Redirect HTTP ke HTTPS
}

server {
    listen 443 ssl http2;
    server_name api.nusanara.app;

    # SSL/TLS via Let's Encrypt (Certbot)
    ssl_certificate     /etc/letsencrypt/live/api.nusanara.app/fullchain.pem;
    ssl_certificate_key /etc/letsencrypt/live/api.nusanara.app/privkey.pem;
    ssl_protocols       TLSv1.2 TLSv1.3;   # Matikan TLS 1.0/1.1 (usang)

    # Konfigurasi CORS untuk frontend Vercel
    add_header Access-Control-Allow-Origin  "https://nusanara.app";
    add_header Access-Control-Allow-Methods "GET, POST, OPTIONS";
    add_header Access-Control-Allow-Headers "Authorization, Content-Type";

    location / {
        proxy_pass          http://127.0.0.1:8000;
        proxy_http_version  1.1;
        proxy_set_header    Host $host;
        proxy_set_header    X-Real-IP $remote_addr;
        proxy_set_header    X-Forwarded-Proto $scheme;

        # KONFIGURASI KRITIS UNTUK SSE STREAMING
        proxy_buffering     off;       # Teruskan token langsung tanpa buffer
        proxy_cache         off;       # Jangan cache respons SSE
        proxy_read_timeout  300s;      # 5 menit untuk respons LLM panjang
        proxy_send_timeout  300s;

        # Header SSE wajib
        add_header          Cache-Control "no-cache";
        add_header          X-Accel-Buffering "no";
    }
}
\end{lstlisting}

\noindent \textbf{Penjelasan konfigurasi Nginx:}
\begin{itemize}
  \item \textbf{HTTPS via Let's Encrypt:} Sertifikat SSL gratis dari Let's Encrypt
    dikelola oleh Certbot yang berjalan sebagai cron job untuk auto-renewal
    setiap 90 hari. HTTPS wajib karena JWT Bearer token dikirim dari frontend ---
    tanpa HTTPS, token dapat dicuri via \textit{man-in-the-middle attack}.
  \item \textbf{\texttt{proxy\_buffering off}:} Konfigurasi paling krusial untuk SSE.
    Nginx secara default meng-\textit{buffer} seluruh respons backend sebelum
    mengirimkannya ke klien. Untuk SSE, ini berarti token yang dihasilkan LLM
    menumpuk di buffer Nginx dan tidak dikirim ke klien hingga respons selesai ---
    menghilangkan efek streaming. Dengan \texttt{proxy\_buffering off}, setiap
    chunk \texttt{data: ...\textbackslash n\textbackslash n} langsung diteruskan.
  \item \textbf{\texttt{X-Accel-Buffering: no}:} Header tambahan yang diakui
    oleh beberapa versi Nginx untuk memastikan buffering benar-benar dinonaktifkan
    bahkan jika ada override konfigurasi upstream.
  \item \textbf{CORS headers:} Membatasi akses hanya dari domain
    \texttt{https://nusanara.app} (frontend Vercel), mencegah akses dari
    domain tidak sah. \texttt{OPTIONS} diperlukan untuk preflight CORS request
    browser.
\end{itemize}

\subsubsection{Konfigurasi PostgreSQL untuk Performa}

PostgreSQL dikonfigurasi melalui \texttt{postgresql.conf} untuk mengoptimalkan
performa pada VPS 16 GB:

\begin{lstlisting}[language=bash, caption={Konfigurasi PostgreSQL Kritis di postgresql.conf}]
# Memori
shared_buffers          = 2GB      # Disesuaikan karena Llama 3.1 8B menyita RAM
effective_cache_size    = 4GB      # Estimasi cache OS yang tersedia setelah Llama
work_mem                = 256MB    # Memori per operasi sort/hash
maintenance_work_mem    = 1GB      # Untuk operasi VACUUM, CREATE INDEX

# Checkpointing
checkpoint_completion_target = 0.9  # Sebar I/O checkpoint
wal_buffers             = 64MB

# Koneksi
max_connections         = 50       # FastAPI menggunakan pool 10 koneksi
\end{lstlisting}

\noindent \textbf{Justifikasi \texttt{shared\_buffers = 2 GB}:}
Nilai 2 GB dialokasikan karena kapasitas RAM (16 GB) sebagian besar telah digunakan
oleh proses Ollama dan sistem operasi (lihat estimasi alokasi pada Tabel~\ref{tab:services}).
\texttt{effective\_cache\_size = 4 GB} memberitahu \textit{planner} PostgreSQL estimasi
total cache yang tersedia agar dapat memilih strategi \textit{query} yang lebih efisien.

\subsubsection{Pipeline Deployment: Dari Kode ke Produksi}

Alur deployment ke VPS dilakukan secara manual menggunakan Git pull, tanpa CI/CD
otomatis. Prosedur deployment standar:

\begin{lstlisting}[language=bash, caption={Prosedur Deployment Manual ke VPS}]
# 1. SSH ke VPS
ssh ubuntu@<ip-vps>

# 2. Pull kode terbaru dari repository
cd /home/ubuntu/nusanara
git pull origin main

# 3. Aktifkan virtual environment dan install dependencies
source .venv/bin/activate
pip install -r backend/requirements.txt --quiet

# 4. Jalankan migrasi database (jika ada perubahan schema)
psql -U nusanara -d nusanara_db -f migrations/latest.sql

# 5. Restart FastAPI service (zero-downtime jika menggunakan 2 workers)
sudo systemctl restart nusanara.service

# 6. Verifikasi status
sudo systemctl status nusanara.service
curl http://127.0.0.1:8000/health | python3 -m json.tool
\end{lstlisting}

\noindent \textbf{Catatan:} Verifikasi health check setelah restart memastikan
koneksi database dan Ollama aktif sebelum traffic diarahkan kembali ke endpoint.

\subsubsection{Keamanan Infrastruktur}

Beberapa konfigurasi keamanan yang diterapkan:
\begin{itemize}[noitemsep]
  \item \textbf{UFW Firewall:} Hanya port 22 (SSH), 80 (HTTP), dan 443 (HTTPS)
    yang terbuka secara publik. Port 5432 (PostgreSQL), 8000 (FastAPI), dan
    11434 (Ollama) terikat ke \texttt{127.0.0.1} (localhost only), tidak dapat
    diakses dari internet.
  \item \textbf{SSH Key Authentication:} Login SSH hanya dengan kunci publik/privat.
    Password authentication dinonaktifkan di \texttt{/etc/ssh/sshd\_config}
    (\texttt{PasswordAuthentication no}).
  \item \textbf{Fail2Ban:} Memblokir IP yang gagal login SSH lebih dari 5 kali
    dalam 10 menit, mencegah \textit{brute-force attack}.
  \item \textbf{Variabel Rahasia:} Seluruh kredensial (DATABASE\_URL, CLERK\_SECRET)
    disimpan di file \texttt{.env} yang di-\textit{gitignore}, dimuat via
    \texttt{EnvironmentFile} systemd.
\end{itemize}
